\subsection*{基本定理的其他证明}
\phantomsection
\addcontentsline{toc}{subsection}{基本定理的其他证明}

在 {[}278 b{]} 中，F. 舒尔证明，在典范坐标中，(15) 的诸 \(\psi_{ik}\) 满足微分方程

\[
\frac {d}{d t} (t \psi_ {i k} (t a)) = \delta_ {i k} + \sum_ {j, l} c _ {j l} ^ {k} t a _ {l} \psi_ {j i} (t a).\tag{25.23}
\]

对这些方程积分，便得到一个等价于公式

\[
\varpi (X) = \sum_ {n \geq 0} \frac {1}{(n + 1) !} (\operatorname{ad} (X)) ^ {n}\tag{25.24}
\]

的公式，它给出指数映射在点 X 处的右微分 \(\varpi(X)\)；特别地，在典范坐标中，诸 \(\psi_{ij}\) 可延拓为诸 \(a_{k}\) 的整函数。F. 舒尔由此导出一个使李先前的一条注记精确化的结果：如果在变换群的定义 (25.4) 中只假设诸 \(f_{i}\) 属于类 \(C^{2}\)，则该群与一个解析群全同构。\(^{12}\) 继他对微分方程组积分的研究之后，E. 嘉当（{[}52 a{]}，v. II₂, p.~371）于 1904 年引入了普法夫型

\[
\omega_ {k} = \sum_ {i = 1} ^ {r} \psi_ {k i} d a _ {i} (1 \leq i \leq r)\tag{25.25}
\]

（用 (25.15) 的记号），后来称为毛雷尔–嘉当形式。于是毛雷尔的条件 (25.17) 可以写成

\[
d \omega_ {k} = - \frac {1}{2} \sum_ {i, j} c _ {i j} ^ {k} \omega_ {i} \wedge \omega_ {j};
\]

E. 嘉当证明，从诸 \(\omega_{k}\) 出发可以展开有限连续群的理论，并确立了这一观点与李的观点的等价性。但对他来说，这个方法的吸引力尤其在于它适用于“无限连续”群——对这些群，他把理论推进得远比利所做的更远——而且它使他得以建立其广义“活动标架”理论。

